\documentclass[letterpaper,11pt]{article}
\usepackage[spanish]{babel}
\usepackage[margin=1in]{geometry}
\usepackage[final]{graphicx}
\usepackage[sc]{mathpazo}
\usepackage{amsmath}
\usepackage{gensymb}
\usepackage{amssymb}
\usepackage{enumerate}
\usepackage{hyperref}
\usepackage{multicol}
\usepackage{caption}
\usepackage{subcaption}
\usepackage{tikz}
\usetikzlibrary{positioning, arrows.meta, shapes, calc}
\usepackage{fancyhdr}

\hypersetup{
    colorlinks=true,
    linkcolor=blue,
    filecolor=magenta,      
    urlcolor=blue,
    pdftitle={Boceto de Arquitectura SUAS},
}

\pagestyle{fancy}
\fancyhf{}
\lfoot{ \fancyplain{}{\texttt{}} }
\rfoot{ \fancyplain{}{\thepage} }
\renewcommand{\headrulewidth}{0pt}
\renewcommand{\footrulewidth}{0pt}
\usepackage{float}

\begin{document}

\begin{center}
\vspace{-0.2cm}
\noindent\makebox[\linewidth]{\rule{\textwidth}{0.4pt}}
\vspace{0.2cm}
\textsc{Proyecto \textit{quadcopter} para SUAS }\\
\vspace{0.1cm}
\huge{\textsc{Boceto de Arquitectura resBot}}\\
\vspace{0.4cm}
\normalsize{\textsc{Vicente Andrade, Antonio Avendaño, Cristóbal Fuentes, Alejandro Mena}}
\vspace{-0.2cm}\\
\noindent\makebox[\linewidth]{\rule{\textwidth}{0.4pt}}
\end{center}

\vspace{1cm}

\section*{Diagrama de Arquitectura}

\vspace{0.5cm}

\begin{center}
\begin{tikzpicture}[
    node distance=1.8cm and 2.5cm,
    block/.style={
        rectangle, draw, thick,
        minimum width=3cm, minimum height=1.5cm,
        align=center, font=\sffamily, fill=gray!10
    },
    line/.style={draw, thick},
    arrow/.style={draw, thick, -{Latex[length=2.5mm]}},
    darrow/.style={draw, thick, {Latex[length=2.5mm]}-{Latex[length=2.5mm]}}
]

% Nodos principales
\node[block] (pms) {\hyperlink{sec:pms}{Power Managment Supply}};
\node[block, below=2cm of pms] (telco) {\hyperlink{sec:telco}{MCU Telecomunicaciones}};
\node[block, right=2.5cm of telco] (raxa) {\hyperlink{sec:radxa}{Radxa Zero 3w}};
\node[block, left=2.5cm of telco] (sec) {\hyperlink{sec:control}{MCU Control Secundario}};

% Nodos inferiores
\node[block, below=1.8cm of raxa] (gimbal) {\hyperlink{sec:camara}{Gimbal / Cámara}};
\node[block, below=1.8cm of telco] (lora) {\hyperlink{sec:lora}{Módulo LoRa}};

% Múltiples módulos FOC (efecto de apilamiento para representar 4 módulos)
\node[block, below=1.8cm of sec, xshift=-0.3cm, yshift=-0.3cm] (foc4) {};
\node[block, below=1.8cm of sec, xshift=-0.2cm, yshift=-0.2cm] (foc3) {};
\node[block, below=1.8cm of sec, xshift=-0.1cm, yshift=-0.1cm] (foc2) {};
\node[block, below=1.8cm of sec] (foc1) {\hyperlink{sec:foc}{FOC}};

% --- Conexiones ---

% Conexiones de Alimentación desde PMS
% A MCU Telco (Centro)
\draw[arrow] (pms) -- node[right] {3V3} (telco);

% A Raxa (5V) (Derecha)
\draw[arrow] ([xshift=0.5cm]pms.south) -- ++(0,-0.7) -| node[pos=0.25, above] {5V} (raxa.north);

% A MCU Control Secundario (3V3) (Izquierda)
\draw[arrow] ([xshift=-0.5cm]pms.south) -- ++(0,-0.7) -| node[pos=0.25, above] {3V3} (sec.north);

% Conexiones SPI (Comunicación)
% Telco <-> Raxa
\draw[darrow] (telco) -- node[above] {SPI} (raxa);

% Telco <-> MCU Control Secundario
\draw[darrow] (sec) -- node[above] {SPI} (telco);

% Telco <-> LoRa
\draw[darrow] (telco) -- node[right] {SPI} (lora);

% Símbolo de antena para LoRa
\draw[thick] (lora.south east) -- ++(0.4, -0.4) coordinate (ant);
\draw[thick] (ant) -- ++(0, 0.5) coordinate (ant_top);
\draw[thick] (ant_top) -- ++(-0.2, 0.2);
\draw[thick] (ant_top) -- ++(0.2, 0.2);
\draw[thick] (ant_top) -- ++(0, 0.25);

% Conexión Raxa a Gimbal/Cámara
\draw[darrow] (raxa) -- node[right] {MIPI CSI-2} (gimbal);

% Conexión MCU Secundario a FOC (4 buses SPI)
\draw[arrow] (sec) -- coordinate (midfoc) (foc1.north);
% Símbolo de bus (diagonal)
\draw[thick] ($(midfoc)+(-0.15,-0.15)$) -- ($(midfoc)+(0.15,0.15)$);
\node[left=0.15cm of midfoc, font=\sffamily] {4};
\node[right=0.15cm of midfoc, font=\sffamily] {SPI};

\end{tikzpicture}
\end{center}

\newpage

% Secciones de destino (hypertargets)
\section{Power Management Supply} \hypertarget{sec:pms}{}
\begin{itemize}
    \item \textbf{Queda diseñar el módulo de alimentación} (Baterías para alimentación independiente).
    \item Conexiones: Entrega 5 [V] a la Radxa Zero 3W y 3.3 [V] a la MCU Control Secundario y a la MCU Telecomunicaciones.
\end{itemize}

\section{MCU de Telecomunicaciones} \hypertarget{sec:telco}{}
\begin{itemize}
    \item Se propone una ESP32 (\textbf{Queda revisar el modelo}).
    \item Conexiones: SPI a Radxa, MCU Control Secundario y Módulo LoRa.
\end{itemize}

\section{Radxa Zero 3W} \hypertarget{sec:radxa}{}
\begin{itemize}
    \item Dedicado principalmente al procesamiento digital de imágenes para cámara.
    \item Posee un procesador de cuatro núcleos (SoC) Rockchip RK3566 de bajo consumo que a la vez posee una NPU de 0.8 TOPS integrada.
    \item \textbf{*Boceto de funcionamiento:} Ejecutar YOLO convertido al formato Rockchip (RKNN) y utilizar OpenCV solo para capturar fotogramas y dibujar las \textit{bounding boxes}.
    \item \textbf{Queda revisar: 4-8GB de RAM (LPDDR4)} para cargar SO, capturar video HD y mantener los tensores YOLO en memoria.
    \item Conexiones: SPI al Módulo de Telecomunicaciones y MIPI CSI-2 a la cámara mediante FPC.
\end{itemize}

\section{MCU de Control Secundario} \hypertarget{sec:control}{}
\begin{itemize}
    \item \textbf{El desarrollo de este módulo será el enfoque del ramo ELO301}
    \item \textbf{Queda revisar el MCU que se utilizará.}
    \item Conexiones: SPI a 4 Módulos FOC y a Módulo Telecomunicaciones.
\end{itemize}

\section{Módulo LoRa} \hypertarget{sec:lora}{}
\begin{itemize}
    \item Se propone un circuito integrado Semtech SX1262 a $\sim$ 915 [MHz] (Alcance de 15-20 [km]) para recibir comandos remotos y transmitir \textbf{datos}.
    \item \textbf{Queda revisar módulos comerciales (Ebyte, Hope RF, Waveshare).}
    \item Conexiones: SPI a Módulo Telecomunicaciones.
\end{itemize}

\section{Gimbal y Cámara} \hypertarget{sec:camara}{}
\begin{itemize}
    \item \textbf{Queda revisar el sensor a utilizar (Ej: IMX519-16MP, IMX477-12MP), el módulo que posea este sensor y si este posee un  conector FPC con separación entre pines de 0.5 [mm] para Radxa. Revisar si se optará por utilizar un lente teleobjetivo (Zoom óptico).} 
\end{itemize}

\section{Módulos FOC} \hypertarget{sec:foc}{}
\begin{itemize}
    \item \textbf{Queda revisar diseño e implementación de módulos FOC.}
    \item \textbf{Queda revisar si se utilizará un Gimbal finalmente y si es el caso, cuál se utilizará.}
\end{itemize}

\newpage
\section*{Enlaces a Componentes}
\begin{itemize}
    \item \href{https://radxa.com/products/zeros/zero3w/#techspec}{Radxa Zero 3w}, Single Board Computer
    \item \href{https://www.semtech.com/products/wireless-rf/lora-connect/sx1262}{SX1262}, LoRa Transceiver
\end{itemize}

\section*{Componentes a evaluar:}
\subsection*{Cámara}
    \begin{itemize}
        \item \href{https://www.arducam.com/imx519-autofocus-camera-module-for-raspberry-pi-arducam-b0371.html}{IMX519-16MP} de Arducam
        \item \href{https://docs.arducam.com/Raspberry-Pi-Camera/Native-camera/12MP-IMX477/#supported-operating-systems}{IMX477-12MP} de Arducam
    \end{itemize}

\newpage
\section*{Glosario}
\begin{itemize}
    \item \textbf{CSI}: Camera Serial Interface
    \item \textbf{FOC}: Field Oriented Control
    \item \textbf{FPC}: Flexible Printed Circuit
    \item \textbf{MCU}: MicroController Unit
    \item \textbf{NPU}: Neural Processing Unit
    \item \textbf{OS}: Operative System
    \item \textbf{PMS}: Power Management Supply 
    \item \textbf{RKNN}: Rockchip Neural Network
    \item \textbf{SoC}: System On a Chip
    \item \textbf{SPI}: Serial Peripheral Interface
    \item \textbf{TOPS}: Tera Operations per second
    \item \textbf{YOLO}: \textit{You Only Look Once} (Algoritmo de Visión Computacional)
\end{itemize}

\end{document}
